\documentclass[11pt]{article}
\usepackage{mathblatt}
% Lösungen nach mathe-nachhilfe gemeinsam.md „Lösungen“ und pruefung.md § 6; Stufenköpfe wortgleich wie im Blatt (Nachzug 12.10.).
\newcommand{\kennung}{GYW}
\newcommand{\pkt}{\textperiodcentered{}}
\begin{document}
\pfheftstil
\pflaufkopf[P10 \pkt{} Prüfung]{sin, cos, tan \pkt{} Lösungen}{\kennung}\pfseitevon
\pfheftkopf{sin, cos, tan}{Lösungen}
\pfstufenzeile{Gleichung aufstellen \pkt{} Winkel gesucht \pkt{} Seite gesucht \pkt{} Pythagoras oder sin, cos, tan?}

\pfstufekopf{Gleichung aufstellen}{letzte 5 Jahre: 2025}
\begin{pfloesungb}{}
\lzb{1.}{a)}{$t$ Hypotenuse, $r$ Gegenkathete, $s$ Ankathete}{$t$ liegt dem rechten Winkel gegenüber; $r$ liegt $\alpha$ gegenüber; $s$ liegt an $\alpha$}
\lzb{}{b)}{$\sin\alpha = \frac{r}{t}$; \ $\cos\alpha = \frac{s}{t}$; \ $\tan\alpha = \frac{r}{s}$}{Gegenkathete durch Hypotenuse; Ankathete durch Hypotenuse; Gegenkathete durch Ankathete}
\end{pfloesungb}
\begin{pfloesungb}{}
\lzb{2.}{a)}{$\alpha$}{$r$ liegt $\alpha$ gegenüber}
\lzb{}{b)}{$\beta$}{$r$ liegt an $\beta$}
\lzb{}{c)}{$\beta$}{$s$ liegt $\beta$ gegenüber, $r$ an $\beta$}
\end{pfloesungb}
\begin{pfloesungb}{}
\lzb{3.}{}{$\sin\gamma = \frac{u}{w}$}{$u$ liegt $\gamma$ gegenüber, $w$ ist die Hypotenuse}
\end{pfloesungb}

\pfstufekopf{Winkel gesucht}{letzte 5 Jahre: 2026 \pkt{} 2024 \pkt{} 2022}
\begin{pfloesungb}{}
\lzb{4.}{a)}{sin}{Gegenkathete und Hypotenuse}
\lzb{}{b)}{cos}{Ankathete und Hypotenuse}
\lzb{}{c)}{tan}{Gegenkathete und Ankathete}
\end{pfloesungb}
\begin{pfloesungb}{}
\lzb{5.}{a)}{$\alpha \approx 66{,}0^\circ$}{$\tan\alpha = \frac{9}{4}$}
\lzb{}{b)}{$\alpha \approx 61{,}3^\circ$}{$\cos\alpha = \frac{2{,}4}{5} = 0{,}48$}
\lzb{}{c)}{$\alpha \approx 18{,}7^\circ$}{$\sin\alpha = \frac{0{,}8}{2{,}5} = 0{,}32$}
\end{pfloesungb}
\begin{pfloesungb}{}
\lzb{6.}{}{$\beta_1 \approx 48{,}6^\circ$}{$\cos\beta_1 = \frac{262}{396} \approx 0{,}662$}
\end{pfloesungb}
\begin{pfloesungb}{}
\lzb{7.}{}{$\varepsilon \approx 65{,}1^\circ$, also etwa $65^\circ$}{Dreieck $BFC$: $\cos\varepsilon = \frac{16}{38} \approx 0{,}421$}
\end{pfloesungb}

\pfstufekopf{Seite gesucht}{letzte 5 Jahre: 2026 \pkt{} 2023 \pkt{} 2022}
\begin{pfloesungb}{}
\lzb{8.}{a)}{A}{$x$ ist Gegenkathete: $\sin 37^\circ = \frac{x}{10}$, $x$ steht oben $\Rightarrow$ mal}
\lzb{}{b)}{B}{$x$ ist Hypotenuse: $\sin 52^\circ = \frac{5}{x}$, $x$ steht unten $\Rightarrow$ geteilt}
\end{pfloesungb}
\begin{pfloesungb}{}
\lzb{9.}{a)}{$x \approx 4{,}1$ cm}{$\sin 31^\circ = \frac{x}{8}$ \ $\Rightarrow$ \ $x = 8 \cdot \sin 31^\circ$}
\lzb{}{b)}{$x \approx 8{,}6$ cm}{$\sin 44^\circ = \frac{6}{x}$ \ $\Rightarrow$ \ $x = 6 : \sin 44^\circ$}
\end{pfloesungb}
\begin{pfloesungb}{}
\lzb{10.}{}{$a \approx 9{,}7$ m}{Dreieck $DBC$: $\sin 49^\circ = \frac{7{,}3}{a}$ \ $\Rightarrow$ \ $a = 7{,}3 : \sin 49^\circ \approx 9{,}67$}
\end{pfloesungb}
\begin{pfloesungb}{}
\lzb{11.}{}{$\overline{BF} = \overline{AF} \approx 91{,}1$ cm; \ $\overline{BC} \approx 172{,}0$ cm}{Dreieck $ABF$: $\overline{BF} = 128{,}9 \cdot \sin 45^\circ \approx 91{,}15$, gleichschenklig; Dreieck $BCF$: $\overline{BC} = \overline{BF} : \cos 58^\circ$}
\end{pfloesungb}

\pfstufekopf{Pythagoras oder sin, cos, tan?}{letzte 5 Jahre: 2026 \pkt{} 2025 \pkt{} 2024 \pkt{} 2022}
\begin{pfloesungb}{}
\lzb{12.}{a)}{Pythagoras}{zwei Seiten gegeben, Seite gesucht, kein Winkel}
\lzb{}{b)}{sin, cos, tan}{Winkel im Spiel}
\lzb{}{c)}{sin, cos, tan}{Winkel gesucht}
\lzb{}{d)}{Pythagoras}{nur Seiten}
\lzb{}{e)}{Winkelsumme}{zwei Winkel gegeben, Winkel gesucht}
\lzb{}{f)}{sin, cos, tan}{Winkel im Spiel}
\end{pfloesungb}
\begin{pfloesungb}{}
\lzb{13.}{}{$x \approx 200{,}7$ cm; \ $\alpha \approx 4{,}9^\circ$}{Pythagoras: $x^2 = 200^2 + 17^2 = 40\,289$; Winkel: $\tan\alpha = \frac{17}{200}$}
\end{pfloesungb}
\end{document}
